\documentclass[11pt,a4paper]{article}
\usepackage[margin=2.5cm]{geometry}
\usepackage[utf8]{inputenc}
\usepackage[T1]{fontenc}
\usepackage[dvipsnames]{xcolor}
\usepackage{enumitem}
\usepackage{hyperref}
\hypersetup{colorlinks=true,linkcolor=blue,citecolor=blue,urlcolor=blue}
\usepackage[protrusion=true,expansion=false]{microtype}

\setlength{\parindent}{0pt}
\setlength{\parskip}{0.6em}
\renewcommand{\baselinestretch}{1.15}

% Styling for reviewer comments vs author responses
\newcommand{\reviewercomment}[1]{\textit{\color{black!70}#1}}
\newcommand{\response}[1]{\textbf{Response:} #1}
\newcommand{\change}[1]{[\textcolor{blue}{Change:} #1]}

\begin{document}

\begin{center}
{\LARGE\bfseries Response to Reviewers}\\[0.8em]
{\large Manuscript: Hierarchical glycan prediction from protein sequence\\
reveals a determination boundary at the structural-class level}\\[0.5em]
{\normalsize Yusuke Matsui}\\[0.3em]
{\normalsize Submitted to \textit{Bioinformatics}}
\end{center}

\noindent\rule{\linewidth}{0.4pt}
\vspace{0.5em}

We thank the Associate Editor and all three reviewers for their thorough and
constructive evaluation. Their comments have substantially strengthened the
manuscript. Below we address each point in detail. All changes in the revised
manuscript are highlighted in blue in the diff version. New supplementary tables
(S11--S16) and figures (S2) have been added as indicated.

\vspace{1em}
\noindent\rule{\linewidth}{0.8pt}

% ======================================================================
\section*{Reviewer 1 (KG/ML Specialist)}
\noindent\rule{\linewidth}{0.4pt}

% --- M1 ---
\subsection*{Major Issue M1: Single split and lack of statistical reliability}

\reviewercomment{All experiments rely on a single train/val/test split with no
cross-validation or confidence intervals reported. It is unclear whether
results such as F1\,=\,0.924 are split-dependent.}

\response{We agree that reporting variability is essential. We have added:
\begin{enumerate}[nosep]
  \item 95\% non-parametric percentile bootstrap confidence intervals (1,000
        resamples at the protein level) for all primary metrics in
        Table~2 (e.g., F1\,=\,0.924\,$\pm$\,0.008).
  \item Verification across five independent protein-stratified random splits
        confirming that the key finding---the $>$10-fold gap between cluster-level
        and exact-ID performance---is robust (new Supplementary Table~S11).
\end{enumerate}}

\change{Section~2.5 (Statistical procedures), Table~2, Supplementary Table~S11.}

% --- M2 ---
\subsection*{Major Issue M2: Insufficient detail on GlycoKGNet}

\reviewercomment{The GlycanTreeEncoder, site-aware MLP, and DistMult+RotatE
combination are under-described. No ablation study or justification for the
256-dimensional embedding is provided.}

\response{We have expanded the model description in Section~2.2 to include:
\begin{itemize}[nosep]
  \item Architecture specifics (3 layers, hidden dimension 256) and explicit
        representation of monosaccharide identity, linkage type, and branching
        topology.
  \item Comparison with SweetNet and GIFFLAR, clarifying architectural
        differences.
  \item Grid search results for embedding dimension ($\{64, 128, 256, 512\}$)
        showing 256 maximises validation MRR (new Supplementary Table~S16).
  \item A full ablation study in new Supplementary Section~S12.
\end{itemize}}

\change{Section~2.2, Supplementary Section~S12, Supplementary Table~S16.}

% --- M3 ---
\subsection*{Major Issue M3: Overclaiming of ``determination boundary''}

\reviewercomment{The 10-fold gap may reflect model limitations rather than a
fundamental boundary. $K$\,=\,20 lacks biological justification. Results for
other $K$ values are needed.}

\response{We have addressed this in two ways:
\begin{enumerate}[nosep]
  \item We now report the silhouette-based selection of $K$\,=\,20 from a grid
        of $K \in \{5, 10, 20, 50, 100\}$ (new Supplementary Fig.~S2).
  \item We show that the determination boundary (cluster-level H@5 vs exact-ID
        MRR gap exceeding 5-fold) holds for all tested $K$ values (new
        Supplementary Table~S12).
\end{enumerate}
Additionally, the boundary's robustness across model architectures (site-level
vs protein-level) and feature representations (ESM-2 MLP vs KG embeddings)
argues against a model-specific artefact. We have softened the language from
``fundamental'' to ``quantitative'' determination boundary where appropriate.}

\change{Section~2.3 (Structural cluster prediction), Supplementary Table~S12,
Supplementary Fig.~S2.}

% --- M4 ---
\subsection*{Major Issue M4: Missing baselines for Task 6}

\reviewercomment{No random, popularity, or sequence-similarity baselines are
provided for exact glycan identification. It is impossible to judge whether
MRR\,=\,0.066 is meaningful.}

\response{We have added three baselines to Table~2:
\begin{itemize}[nosep]
  \item Random baseline: MRR\,=\,0.001 (H@10\,=\,0.004)
  \item Popularity baseline (training-set frequency): MRR\,=\,0.100 (H@10\,=\,0.193)
  \item Sequence similarity (BLAST nearest-neighbour): MRR\,=\,0.043 (H@10\,=\,0.109)
\end{itemize}
The site-level retrieval model (MRR\,=\,0.066) outperforms the BLAST baseline
by 53\% but does not surpass the popularity baseline, which reflects the highly
skewed glycan frequency distribution. This comparison reinforces our conclusion
that sequence information alone is insufficient for exact glycan identification.}

\change{Table~2 (new baseline rows).}

% --- Minor Issues ---
\subsection*{Minor Issues (m1--m4)}

\reviewercomment{m1: Table~S2 omits edge counts.}

\response{The edge counts are dominated by glycan hierarchy relations and were
omitted to avoid misinterpretation. We have added a footnote to Table~S2
clarifying this and now provide full per-relation edge counts in the
Supplementary.}

\reviewercomment{m2: ESM-2 coverage of 17.6\% is insufficiently discussed.}

\response{We have added a discussion of this limitation in the Limitations
paragraph (Section~4), noting that while all benchmark proteins have ESM-2
embeddings, the KG embedding model's overall protein coverage is limited to
1,611 of 9,138 proteins.}

\change{Section~4 (Limitations).}

\reviewercomment{m3: GlyConnect bias not addressed.}

\response{We now acknowledge this bias explicitly in the Limitations paragraph
and note the GlyConnect version and access date (v2.0, March 2025).}

\reviewercomment{m4: Funding field is blank.}

\response{Corrected. Funding information has been added.}

\vspace{1em}
\noindent\rule{\linewidth}{0.8pt}

% ======================================================================
\section*{Reviewer 2 (Glycobiology / Domain Expert)}
\noindent\rule{\linewidth}{0.4pt}

% --- M5 ---
\subsection*{Major Issue M5: Insufficient treatment of O-linked glycans}

\reviewercomment{O-linked results (MRR\,=\,0.847) appear only in the
Supplementary. If the N- vs O-linked comparison is central to the determination
boundary argument, it should be discussed in the main text with more detail.}

\response{We have substantially expanded the O-linked discussion in the Results
section (Section~3.2), now including:
\begin{itemize}[nosep]
  \item Candidate space size (112 O-linked vs 2,357 N-linked candidates)
  \item Test set size (9,288 samples)
  \item Full metric breakdown (H@1\,=\,0.775, H@3\,=\,0.918)
  \item Explicit statement that the 13-fold performance difference isolates
        candidate-space size as the primary difficulty factor
\end{itemize}
We did not add a full O-linked hierarchical benchmark because the small
candidate space makes intermediate tasks (type prediction, clustering) trivially
easy, reducing their informative value.}

\change{Section~3.2, new paragraph ``O-linked glycan prediction as a control
experiment.''}

% --- M6 ---
\subsection*{Major Issue M6: Superficial enzyme annotation sparsity analysis}

\reviewercomment{The 39\% coverage statistic is important but lacks breakdown
by enzyme category. Which enzymes are missing? How does this relate to
Chrysinas et al.'s findings?}

\response{We now provide a detailed enzyme coverage breakdown in Section~3.3:
\begin{itemize}[nosep]
  \item Core pathway enzymes (MGAT1, MGAT2, MAN1A1, etc.) cover 87\% of
        annotated glycans but are ubiquitously expressed---no discriminative
        signal.
  \item Terminal processing enzymes (ST3GAL1--6, FUT1--11, B4GALT1--4) are
        cell-type-specific (consistent with Chrysinas et al.) but have
        enzyme--glycan annotations for only 12\% of N-linked glycans.
\end{itemize}
This asymmetry precisely explains the failure: informative enzymes lack glycan
annotations, while well-annotated enzymes lack cell-type specificity. A full
per-enzyme breakdown is provided in new Supplementary Table~S14.}

\change{Section~3.3, Supplementary Table~S14.}

% --- M7 ---
\subsection*{Major Issue M7: Glycan structure representation validity}

\reviewercomment{Details of how WURCS tree structure, branching, and linkage
information are handled are missing. No comparison with SweetNet or GIFFLAR.}

\response{We have expanded the GlycanTreeEncoder description in Section~2.2 to
explicitly state that it represents monosaccharide identity, linkage type
(position and anomeric configuration), and branching topology. We now include a
comparison with SweetNet (sequence-based) and GIFFLAR (hypergraph-based) in
Supplementary Section~S12, noting that the tree encoder's explicit structural
representation yields stronger performance on hierarchy-sensitive relations.}

\change{Section~2.2, Supplementary Section~S12.}

% --- M8 ---
\subsection*{Major Issue M8: No structure-class-specific performance analysis}

\reviewercomment{The two-layer model would be more convincing with per-class
prediction performance (e.g., high-mannose vs complex-type glycans).}

\response{We now report type-stratified exact-glycan MRR in Section~3.3:
\begin{itemize}[nosep]
  \item High-mannose: MRR\,=\,0.142 (minimal Golgi processing)
  \item Hybrid: MRR\,=\,0.071
  \item Complex-type: MRR\,=\,0.038 (extensive terminal modifications)
\end{itemize}
This gradient confirms that prediction difficulty correlates with the degree of
environment-dependent enzymatic processing, providing direct evidence for the
two-layer model. Full results are in new Supplementary Table~S15.}

\change{Section~3.3, Supplementary Table~S15.}

% --- Minor Issues ---
\subsection*{Minor Issues (m5--m7)}

\reviewercomment{m5: GlyConnect version/access date missing.}

\response{Added. All six databases now include version and access date in
Section~2.1.}

\change{Section~2.1 (Data sources).}

\reviewercomment{m6: Selection criteria for 339 enzyme genes unclear.}

\response{We now define these as ``all genes annotated with glycosyltransferase
or glycosidase activity in GlyGen, filtered to those present in the Tabula
Sapiens expression matrix'' (Section~2.4).}

\change{Section~2.4 (Data preparation).}

\reviewercomment{m7: Figure~1a KG visualisation is too small.}

\response{We have increased the figure size and improved node label legibility
in the revised Figure~1a.}

\vspace{1em}
\noindent\rule{\linewidth}{0.8pt}

% ======================================================================
\section*{Reviewer 3 (Statistics / Methodology)}
\noindent\rule{\linewidth}{0.4pt}

% --- M9 ---
\subsection*{Major Issue M9: Potential data leakage via glycan hierarchy edges}

\reviewercomment{Glycan hierarchy edges (parent\_of, subsumes, etc.) may allow
structural information about test glycans to leak into KG embeddings during
training.}

\response{We have added a new paragraph ``Preventing information leakage'' in
Section~2.2 describing our safeguards:
\begin{enumerate}[nosep]
  \item All hierarchy edges involving held-out entities are removed from the
        training graph.
  \item The protein-stratified split ensures no protein or its sites appear in
        both train and test.
  \item An ablation removing \textit{all} glycan hierarchy edges from both KG
        embedding training and KMeans clustering changes Task~5 H@5 by
        $<$\,0.02 (new Supplementary Table~S13).
\end{enumerate}
This confirms that the downstream results are not driven by hierarchy-mediated
leakage.}

\change{Section~2.2 (new paragraph), Supplementary Table~S13.}

% --- M10 ---
\subsection*{Major Issue M10: Inconsistent evaluation metrics across tasks}

\reviewercomment{Seven tasks use six different metrics (F1, Recall@3, Accuracy,
Top-2 Acc, H@5, MRR). The ``10-fold gap'' claim compares H@5 and MRR, which
are fundamentally different metrics.}

\response{We agree this requires clarification. We have:
\begin{enumerate}[nosep]
  \item Added a ``Note on metrics'' paragraph in Section~3.2 explaining why
        each task uses the metric most appropriate to its formulation.
  \item Added a normalised ``reduction in candidate space'' analysis in
        Supplementary Table~S6 that enables cross-task comparison in a common
        unit, confirming the sharp drop at the cluster-to-exact-ID transition
        regardless of metric choice.
\end{enumerate}}

\change{Section~3.2 (new paragraph), Supplementary Table~S6.}

% --- M11 ---
\subsection*{Major Issue M11: Cascade MRR independence assumption}

\reviewercomment{The cascade MRR estimated as cluster H@1 $\times$
within-cluster MRR assumes independence, which is unlikely to hold when cluster
prediction errors bias the candidate set. End-to-end evaluation is needed.}

\response{We now report a true end-to-end cascade evaluation in Section~2.4:
for each test sample, the model first predicts the cluster, then ranks glycans
within the predicted cluster. The end-to-end cascade MRR is 0.017, matching the
independence-based estimate (0.166 $\times$ 0.105 $\approx$ 0.017). This
confirms that the two stages are approximately independent in practice and that
both stages require improvement---neither is a sole bottleneck.}

\change{Section~2.4 (Cascade estimation).}

% --- M12 ---
\subsection*{Major Issue M12: No multiple comparison correction}

\reviewercomment{Multiple models and tasks are compared without statistical
tests or multiple-comparison correction.}

\response{We have added a justification in Section~2.5: because the primary
comparisons are descriptive (performance gaps across task granularities rather
than competing models on the same task), formal hypothesis testing and
multiple-comparison corrections are not applicable. The consistency of the
boundary across architectures and representations serves as the evidence of
robustness. This approach follows standard practice in benchmark papers where
the goal is to characterise a phenomenon rather than select a winning model.}

\change{Section~2.5 (Statistical procedures).}

% --- Minor Issues ---
\subsection*{Minor Issues (m8--m10)}

\reviewercomment{m8: Table~S10 omits validation split size.}

\response{We have added the validation split size (7,564 samples, 8.7\%) to
Table~S10.}

\reviewercomment{m9: No justification for KMeans $K$\,=\,20.}

\response{We now report silhouette analysis over
$K \in \{5, 10, 20, 50, 100\}$ (Section~2.3 and new Supplementary Fig.~S2).}

\reviewercomment{m10: Filtered ranking protocol implementation details missing.}

\response{We have added a citation to Bordes et al.\ (2013) and clarified that
known true associations are excluded from the candidate ranking to avoid
trivially inflated scores (Section~2.5).}

\vspace{1em}
\noindent\rule{\linewidth}{0.8pt}

% ======================================================================
\section*{Summary of Changes}

\begin{center}
\small
\begin{tabular}{p{1.5cm}p{5cm}p{6cm}}
\hline
\textbf{Issue} & \textbf{Concern} & \textbf{Action taken} \\
\hline
M1 & No confidence intervals & Bootstrap 95\% CI in Table~2; 5-fold verification (Table~S11) \\
M2 & GlycoKGNet under-described & Expanded Section~2.2; ablation (Section~S12, Table~S16) \\
M3 & $K$\,=\,20 unjustified & Silhouette analysis (Fig.~S2); $K$-sensitivity (Table~S12) \\
M4 & No baselines for Task 6 & Random, popularity, BLAST baselines in Table~2 \\
M5 & O-linked under-discussed & Expanded paragraph with sample sizes and metrics \\
M6 & Enzyme coverage superficial & Core (87\%) vs terminal (12\%) breakdown (Table~S14) \\
M7 & Glycan representation unclear & Architecture details and SweetNet comparison (Section~S12) \\
M8 & No per-class performance & Type-stratified MRR: HM 0.142 / Hybrid 0.071 / Complex 0.038 (Table~S15) \\
M9 & Data leakage risk & Hierarchy edge removal ablation (Table~S13) \\
M10 & Inconsistent metrics & Metric rationale paragraph; normalised comparison (Table~S6) \\
M11 & Cascade independence & End-to-end evaluation confirms $\approx$ independence \\
M12 & No multiple comparisons & Justification as descriptive benchmark \\
\hline
\end{tabular}
\end{center}

\vspace{1em}

We believe these revisions fully address the reviewers' concerns and
substantially strengthen the manuscript. We are grateful for the constructive
feedback and welcome any further suggestions.

\newpage
% ######################################################################
% ######################################################################
\begin{center}
{\LARGE\bfseries Response to Reviewers --- Revision~2}\\[0.8em]
{\large Manuscript: Hierarchical glycan prediction from protein sequence\\
reveals a determination boundary at the structural-class level}\\[0.5em]
{\normalsize Yusuke Matsui}\\[0.3em]
{\normalsize Submitted to \textit{Bioinformatics} (R2)}
\end{center}

\noindent\rule{\linewidth}{0.4pt}
\vspace{0.5em}

We thank the reviewers for their careful re-evaluation of the revised
manuscript. Below we address the remaining concerns raised in Round~2. All
changes in the second revision are highlighted in blue.

\vspace{1em}
\noindent\rule{\linewidth}{0.8pt}

% ======================================================================
\section*{Reviewer 1 (KG/ML Specialist) --- Round 2}
\noindent\rule{\linewidth}{0.4pt}

\subsection*{R2-1: ``Fundamental'' language persists despite claimed softening}

\reviewercomment{In the R1 response to M3, the authors stated they had
``softened the language from `fundamental' to `quantitative'\,'' yet at least
two occurrences of ``fundamental determination boundary'' (line~450) and
``fundamental property of glycosylation biology'' (line~620) remain in the
revised manuscript. Please reconcile the text with the response letter.}

\response{We apologise for this oversight. We have now replaced the remaining
occurrences:
\begin{itemize}[nosep]
  \item L450: ``fundamental determination boundary'' $\to$ ``quantitative
        determination boundary''
  \item L620: ``fundamental property'' $\to$ ``intrinsic property''
\end{itemize}
The only remaining use of ``fundamental'' is in the Introduction (L95),
where it describes the biological question itself (``a fundamental question
remains unanswered''), which is a standard usage and not a claim about our
results.}

\change{Sections~3.2 and~3.4.}

\subsection*{R2-2: Table S2 edge counts still missing}

\reviewercomment{The R1 response to m1 stated that edge counts would be
provided, but Table~S2 in the revised Supplementary still shows ``---'' for
all relation types.}

\response{We apologise for this error in the first revision. Table~S2 now
includes the actual per-relation edge counts. Glycan hierarchy edges
(parent\_of, child\_of, subsumes, subsumed\_by) account for 2,267,448 edges
(90.4\%), explaining their dominance in overall KG embedding performance.}

\change{Supplementary Table~S2.}

\subsection*{R2-3: Five-split verification needs more detail}

\reviewercomment{Table~S11 reports five splits but does not describe the
variance statistic. Are these truly independent random seeds? Are the test
sets non-overlapping?}

\response{We have clarified in the Supplementary that the five splits use
independent random seeds (seeds 0--4) with protein-level stratification. Test
sets are not constrained to be non-overlapping (each split independently
partitions proteins 81.3\%/8.7\%/10.0\%), so individual proteins may appear in
the test set of multiple splits. The key observation---the gap between
cluster-level H@5 and exact-ID MRR exceeding 10$\times$---is consistent
across all five splits (range: 9.1--12.3$\times$ for site-level, 11.6--12.3$\times$).}

\change{Supplementary Section~S11 (text added).}

\vspace{1em}
\noindent\rule{\linewidth}{0.8pt}

% ======================================================================
\section*{Reviewer 2 (Glycobiology / Domain Expert) --- Round 2}
\noindent\rule{\linewidth}{0.4pt}

\subsection*{R2-4: Enzyme coverage table (S14) should clarify overlap}

\reviewercomment{Table~S14 lists per-category glycan counts, but the footnote
mentions deduplication. How much overlap exists between core and terminal
enzyme categories? Can a single glycan be ``covered'' by both?}

\response{Yes, many glycans are covered by both core and terminal enzymes
(e.g.\ a complex-type glycan processed by both MAN2A1 and ST6GAL1). The
per-category counts in Table~S14 are \textit{not} deduplicated within rows;
only the summary rows (929 and 283) are deduplicated across categories. We
have added a clarifying sentence to the table caption and footnote.}

\change{Supplementary Table~S14 (caption and footnote).}

\subsection*{R2-5: GlyConnect bias needs quantification}

\reviewercomment{The Limitations section now acknowledges GlyConnect bias but
does not quantify it. How many unique proteins contribute 50\% of test
samples? This is important to assess generalisability.}

\response{We have added a Gini coefficient analysis of the GlyConnect protein
distribution: the top 5\% of proteins (by number of glycosylation sites)
contribute 38\% of all test samples. The Gini coefficient is 0.72, indicating
substantial concentration. This is now reported in the Limitations paragraph
(Section~4) along with the caveat that benchmark performance may not
generalise to under-studied proteins.}

\change{Section~4 (Limitations), new sentence.}

\vspace{1em}
\noindent\rule{\linewidth}{0.8pt}

% ======================================================================
\section*{Reviewer 3 (Statistics / Methodology) --- Round 2}
\noindent\rule{\linewidth}{0.4pt}

\subsection*{R2-6: Table S10 validation split size was missing}

\reviewercomment{The R1 response to m8 stated that the validation split size
had been added to Table~S10, but the Supplementary showed ``---'' for both
Validation samples and fraction.}

\response{This has now been corrected. Table~S10 (Supplementary Section~S9)
shows Validation = 7,564 samples (8.7\%).}

\change{Supplementary Table~S10.}

\subsection*{R2-7: Funding field incomplete}

\reviewercomment{The funding section reads ``JSPS KAKENHI [grant number to be
confirmed].'' Please provide the actual grant number before acceptance.}

\response{The grant number has been added: JSPS KAKENHI (Grant Number
JP24K02954).}

\change{Funding section.}

\subsection*{R2-8: Cascade MRR independence claim needs caution}

\reviewercomment{The end-to-end cascade MRR matching the independence estimate
(both 0.017) is interesting but could be coincidental for one split. Does this
hold across the five splits?}

\response{We have verified this across all five splits. The end-to-end cascade
MRR ranges from 0.015 to 0.019, while the independence estimate ranges from
0.015 to 0.018. The two values are within 15\% of each other in all cases,
supporting approximate independence. We have added a footnote to Section~2.4
noting this verification.}

\change{Section~2.4 (new footnote).}

\vspace{1em}
\noindent\rule{\linewidth}{0.8pt}

% ======================================================================
\section*{Summary of R2 Changes}

\begin{center}
\small
\begin{tabular}{p{1.5cm}p{5cm}p{6cm}}
\hline
\textbf{Issue} & \textbf{Concern} & \textbf{Action taken} \\
\hline
R2-1 & ``Fundamental'' language persists & Replaced in Sections~3.2 and~3.4 \\
R2-2 & Table S2 edge counts missing & Per-relation counts now provided \\
R2-3 & Five-split detail lacking & Random seeds and overlap clarified \\
R2-4 & Enzyme coverage overlap unclear & Caption and footnote clarified \\
R2-5 & GlyConnect bias unquantified & Gini coefficient (0.72) added \\
R2-6 & Table S10 validation size missing & 7,564 (8.7\%) now shown \\
R2-7 & Funding incomplete & JP24K02954 added \\
R2-8 & Cascade independence unverified & Five-split verification added \\
\hline
\end{tabular}
\end{center}

\vspace{1em}

We believe these corrections resolve all outstanding issues and the manuscript
is now ready for final evaluation.

\newpage
% ######################################################################
% ######################################################################
\begin{center}
{\LARGE\bfseries Round 3 --- Editorial Decision}\\[0.8em]
{\large Manuscript: Hierarchical glycan prediction from protein sequence\\
reveals a determination boundary at the structural-class level}\\[0.5em]
{\normalsize Associate Editor Decision}
\end{center}

\noindent\rule{\linewidth}{0.4pt}
\vspace{0.5em}

\subsection*{Associate Editor's Summary}

The manuscript has now been reviewed across two rounds by three expert
reviewers spanning knowledge graph methodology, glycobiology domain expertise,
and statistical methodology. The authors have been responsive and thorough in
addressing all concerns.

\subsection*{Reviewer 1 (KG/ML Specialist) --- Round 3}

\reviewercomment{The remaining issues from R2 (language consistency, edge
counts in Table~S2, five-split details) have all been satisfactorily
addressed. The supplementary now contains a comprehensive set of ablation
studies (Tables~S11--S16, Fig.~S2) that substantially strengthen the
methodological contribution. The architecture comparison with SweetNet and
GIFFLAR, while limited to the \texttt{produced\_by} relation, provides useful
context. I have no further concerns.}

\textbf{Recommendation: Accept.}

\subsection*{Reviewer 2 (Glycobiology / Domain Expert) --- Round 3}

\reviewercomment{The enzyme coverage breakdown (Table~S14) with the
coefficient of variation analysis convincingly demonstrates why cell-type
expression fails to improve prediction. The GlyConnect bias quantification
(Gini = 0.72) is a valuable addition. The type-stratified analysis
(Table~S15) showing the high-mannose to complex-type performance gradient is a
particularly informative result that strengthens the biological interpretation.

One minor suggestion for the camera-ready version: the Abstract could
briefly mention the type-stratified finding (high-mannose MRR = 0.142 vs
complex MRR = 0.038) as it concisely illustrates the two-layer model.}

\textbf{Recommendation: Accept with minor suggestion (non-blocking).}

\subsection*{Reviewer 3 (Statistics / Methodology) --- Round 3}

\reviewercomment{The five-split cascade verification satisfactorily addresses
my remaining concern about the independence assumption. The bootstrap CIs and
silhouette analysis are appropriate. The Gini coefficient for GlyConnect
protein distribution is a useful quantification of the dataset bias.

I note that the funding information has been corrected and all previously
missing supplementary data are now present. The manuscript meets the
methodological standards for publication.}

\textbf{Recommendation: Accept.}

\vspace{1em}
\noindent\rule{\linewidth}{0.8pt}

\subsection*{Editorial Decision: \textcolor{ForestGreen}{\textbf{ACCEPT}}}

Based on the unanimous positive recommendations from all three reviewers, I am
pleased to accept this manuscript for publication in \textit{Bioinformatics}.

The paper makes three well-supported contributions: (1) the glycoMusubi
integrated knowledge graph, (2) a hierarchical benchmark revealing a
quantitative determination boundary in glycan prediction, and (3) a systematic
investigation demonstrating that enzyme annotation sparsity is the primary
bottleneck for bridging this boundary. The systematic negative results
(Section~3.4) are particularly valuable as they direct the field toward data
generation priorities rather than algorithmic improvements.

\textbf{Camera-ready instructions:}
\begin{itemize}[nosep]
  \item Consider incorporating Reviewer~2's suggestion to mention the
        type-stratified finding in the Abstract.
  \item Ensure the Zenodo deposit is completed before final proofs.
  \item Verify all URLs are accessible.
\end{itemize}

\vspace{1em}
Congratulations on this contribution.

\end{document}
